% Volume II: Causal Explosion
% Part III. The Combinatorics of Context
% Chapter 14: Ecological Judgment
% Spine: proof-spines/ch014-spine.md

\chapter{Ecological Judgment}
\label{ch:ch014}

\textbf{Ecological judgment} begins from the fact that context can change the causal type of an act. What looks local may depend on technical systems, household arrangements, institutional incentives, medicine, psychology, labor, law, or environmental conditions. To formalize that embedding, the chapter introduces a multiscale model
\[
X^{(0)}\subset X^{(1)}\subset\cdots\subset X^{(L)},
\]
\Definition\ This multiscale embedding fixes the chapter's vocabulary for ecological levels of context.
where level \(0\) is the immediately visible action and higher levels represent interpersonal, institutional, infrastructural, historical, and ecological surroundings.

The central problem is to choose the smallest scale at which the event becomes intelligible without letting inquiry expand indefinitely. A rule, heuristic, burden, or presumption should therefore be judged by how well it fits the environment described at the relevant level of \(X^{(L)}\), not by an abstract fantasy of omniscience. Paranoia is rational here only when it is ecologically calibrated: enough context to identify the act correctly, not so much that explanation dissolves into boundless scenery.
